\documentclass[11pt]{article}
\usepackage{amsmath,amssymb,amsthm}
\usepackage{algorithm}
\usepackage[noend]{algpseudocode}
\usepackage{fullpage}
\usepackage{hyperref}
\usepackage{xcolor}
\newtheorem{theorem}{Theorem}
\newtheorem{lemma}{Lemma}
\newtheorem{assumption}{Assumption}
\newcommand{\E}{\mathbb{E}}
\newcommand{\R}{\mathbb{R}}

\begin{document}

\title{Top‑\(K\) Stochastic Gradient Descent\\
(Non‑convex Smooth Setting)}
\author{}
\date{}
\maketitle

%%%%%%%%%%%%%%%%%%%%%%%%%%%%%%%%%%%%%%%%%%%%%%%%%%%%%%%%%%%%
\section*{Algorithm Details Expansion}

\subsection*{Mathematical Formulation}
Consider a stochastic optimisation problem
\[
\min_{w\in\R^{d}} \; f(w)\;:=\;\E_{z\sim\mathcal{D}}[F(w;z)],
\]
where \(F(\cdot;z)\) is differentiable but possibly non‑convex.
At iteration \(t\) a stochastic gradient \(g_t\) is sampled:
\[
g_t \;:=\; \nabla F(w_t;z_t), \qquad z_t\sim\mathcal{D}.
\]

Define the \emph{Top‑\(K\) compression operator}
\[
\mathcal{C}_K:\R^{d}\to\R^{d},\qquad
\big[\mathcal{C}_K(v)\big]_i=
\begin{cases}
v_i & \text{if }i\in\mathcal{I}_K(v),\\[2mm]
0   & \text{otherwise},
\end{cases}
\]
where \(\mathcal{I}_K(v)\) contains the indices of the \(K\) components of
\(v\) with largest absolute value \(|v_i|\) (ties broken arbitrarily).
The sparse gradient used for the update is
\[
\tilde g_t \;:=\; \mathcal{C}_K(g_t).
\]

With learning rate \(\eta_t>0\) the Top‑\(K\) SGD update reads
\[
\boxed{ \; w_{t+1}= w_t - \eta_t \,\tilde g_t \;}
\tag{1}
\]

Hyper‑parameters:
\begin{itemize}
\item learning‑rate schedule \(\{\eta_t\}_{t\ge0}\) (e.g. \(\eta_t=\eta/\sqrt{t+1}\));
\item sparsity level \(K\in\{1,\dots,d\}\) (or equivalently compression ratio
\(\rho:=K/d\in(0,1]\)).
\end{itemize}

\subsection*{Pseudocode}
\begin{algorithm}[H]
\caption{Top‑\(K\) Stochastic Gradient Descent}
\begin{algorithmic}[1]
\State \textbf{Input:} initial point \(w_{0}\in\R^{d}\), learning‑rate schedule
\(\{\eta_t\}\), sparsity \(K\).
\For{$t=0,1,2,\dots,T-1$}
    \State Sample a minibatch (or a single datum) \(z_t\sim\mathcal{D}\).
    \State Compute stochastic gradient \(g_t\gets\nabla F(w_t;z_t)\).
    \State \(\tilde g_t\gets\mathcal{C}_K(g_t)\)   \Comment{keep top‑\(K\) entries}
    \State Update \(w_{t+1}\gets w_t-\eta_t\,\tilde g_t\).
\EndFor
\State \textbf{Return} \(w_T\) (or optionally the averaged iterate
\(\bar w_T:=\frac{1}{T}\sum_{t=0}^{T-1}w_t\)).
\end{algorithmic}
\end{algorithm}

\subsection*{Dry Run Example}
 minimise \(f(w)=(w-3)^2\) (scalar, \(d=1\)).  
 Gradient: \(\nabla f(w)=2(w-3)\).  
 Choose \(w_0=0\), learning rate \(\eta=0.1\), and \(K=1\) (the only coordinate is kept).

\begin{center}
\begin{tabular}{c|c|c|c}
$t$ & $w_t$ & $g_t=2(w_t-3)$ & $w_{t+1}=w_t-\eta g_t$ \\ \hline
0 & $0$ & $-6$ & $0-0.1(-6)=0.6$  \\
1 & $0.6$ & $2(0.6-3)=-4.8$ & $0.6-0.1(-4.8)=1.08$ \\
2 & $1.08$ & $2(1.08-3)=-3.84$ & $1.08-0.1(-3.84)=1.464$ \\
\end{tabular}
\end{center}
Corresponding objective values:
\(f(0)=9,\; f(0.6)=5.76,\; f(1.08)=3.6864,\; f(1.464)=2.334\).
The iterates move towards the minimiser \(w^\star=3\) even though only the
single (top‑\(K\)) coordinate is communicated (trivial in 1‑D).

%%%%%%%%%%%%%%%%%%%%%%%%%%%%%%%%%%%%%%%%%%%%%%%%%%%%%%%%%%%%
\section*{Convergence Theory Development}

\subsection*{Assumptions}
\begin{assumption}[Smoothness]\label{ass:smooth}
\(f\) is \(L\)-smooth, i.e.
\[
\forall u,v\in\R^{d}:\qquad
f(u)\le f(v)+\langle\nabla f(v),u-v\rangle+\frac{L}{2}\|u-v\|^{2}.
\tag{A1}
\]
\end{assumption}

\begin{assumption}[Unbiased Stochastic Gradients]\label{ass:unbiased}
\[
\E[g_t\mid w_t]=\nabla f(w_t),\qquad
\E\big[\|g_t\|^{2}\mid w_t\big]\le G^{2}
\tag{A2}
\]
for some constant \(G>0\).
\end{assumption}

\begin{assumption}[Compression Operator]\label{ass:compression}
The Top‑\(K\) operator \(\mathcal{C}_K\) satisfies (see e.g.~\cite{Alistarh2017})
\[
\E\!\big[\|\mathcal{C}_K(v)-v\|^{2}\big]\le (1-\rho)\,\|v\|^{2},
\qquad
\rho:=\frac{K}{d}\in(0,1].
\tag{A3}
\]
The expectation is taken over the random tie‑breaking (if any) of the
Top‑\(K\) selection.
\end{assumption}

\begin{assumption}[Learning‑rate]\label{ass:lr}
The step size is non‑increasing and satisfies
\[
\sum_{t=0}^{\infty}\eta_t =\infty, \qquad
\sum_{t=0}^{\infty}\eta_t^{2}<\infty .
\tag{A4}
\]
A concrete choice used in the theorem is \(\eta_t=\frac{\eta}{\sqrt{t+1}}\)
with a constant \(\eta>0\).
\end{assumption}

\subsection*{Main Convergence Theorem}
\begin{theorem}[Convergence of Top‑\(K\) SGD]\label{thm:main}
Under Assumptions \ref{ass:smooth}–\ref{ass:lr}, let \(\{w_t\}\) be generated by
\eqref{1}.  Define the averaged gradient norm
\[
\Phi_T \;:=\; \frac{1}{T}\sum_{t=0}^{T-1}
\E\!\big[\|\nabla f(w_t)\|^{2}\big].
\]
Then for any \(T\ge1\)
\[
\Phi_T \;\le\;
\frac{2\big(f(w_0)-f^\star\big)}{\eta\sqrt{T}}
\;+\;
\frac{L G^{2}\big(1+(1-\rho)\big)}{\rho}\,
\frac{\eta}{\sqrt{T}},
\tag{2}
\]
where \(f^\star:=\inf_{w}f(w)\).  Consequently,
\[
\Phi_T = \mathcal{O}\!\Big(\frac{1}{\sqrt{T}}\Big).
\]
\end{theorem}

\subsection*{Theoretical Properties}
\begin{itemize}
\item \textbf{Stability w.r.t. sparsity.}
The bound \eqref{2} degrades gracefully as \(\rho\downarrow 0\);
the additional term scales as \(\frac{1-\rho}{\rho}\), reflecting the variance
introduced by aggressive sparsification.
\item \textbf{Hyper‑parameter sensitivity.}
The optimal constant step size \(\eta\) balances the two terms in \eqref{2} and
is of order \(\sqrt{\frac{2\big(f(w_0)-f^\star\big)}{L G^{2}(1+(1-\rho))/\rho}}\).
\item \textbf{Comparison to vanilla SGD.}
When \(\rho=1\) (i.e. \(K=d\)) the bound reduces to the standard
\(\mathcal{O}(1/\sqrt{T})\) rate for non‑convex SGD.
\end{itemize}

\subsection*{Discussion}
The theorem shows that even with extreme communication reduction (e.g.
\(K\ll d\)), Top‑\(K\) SGD retains the same asymptotic
\(\mathcal{O}(1/\sqrt{T})\) decay of the expected squared gradient norm.
The price is a larger constant factor proportional to \((1-\rho)/\rho\).
In practice this factor is often offset by the substantial reduction
in communication cost, especially for high‑dimensional deep‑learning models.

%%%%%%%%%%%%%%%%%%%%%%%%%%%%%%%%%%%%%%%%%%%%%%%%%%%%%%%%%%%%
\section*{Complete Convergence Proof}
\subsection*{Preliminaries (Definitions and Lemmas)}
\begin{lemma}[Smoothness inequality]\label{lem:smooth}
If \(f\) is \(L\)-smooth then for any \(u,v\in\R^{d}\)
\[
f(u) \le f(v)+\langle\nabla f(v),u-v\rangle+\frac{L}{2}\|u-v\|^{2}.
\]
\end{lemma}
\begin{proof}
Standard; follows from integrating the Lipschitz condition on the gradient.
\end{proof}

\begin{lemma}[Compression variance bound]\label{lem:comp}
For any random vector \(v\) and compression operator \(\mathcal{C}_K\) satisfying
Assumption \ref{ass:compression},
\[
\E\!\big[\|\mathcal{C}_K(v)-v\|^{2}\big]\le (1-\rho)\,\E\!\big[\|v\|^{2}\big].
\]
\end{lemma}
\begin{proof}
Directly the statement of Assumption \ref{ass:compression}.
\end{proof}

\subsection*{Main Proof (Step‑by‑Step Derivation)}
We now prove Theorem \ref{thm:main}.  
All expectations are taken w.r.t. the randomness up to time \(t\)
(including the stochastic gradient and the random tie‑breaking in the Top‑\(K\)
operator).

\bigskip
\noindent\textbf{[Step 1] (Smoothness applied to the update).}
Using Lemma~\ref{lem:smooth} with \(u=w_{t+1}\) and \(v=w_t\):
\[
f(w_{t+1})\le f(w_t)
+\langle\nabla f(w_t),w_{t+1}-w_t\rangle
+\frac{L}{2}\|w_{t+1}-w_t\|^{2}.
\tag{3}
\]

\noindent\textbf{[Step 2] (Insert the update rule).}
From \eqref{1}, \(w_{t+1}-w_t=-\eta_t\tilde g_t\).  Substituting into \eqref{3}:
\[
f(w_{t+1})\le f(w_t)
-\eta_t\langle\nabla f(w_t),\tilde g_t\rangle
+\frac{L\eta_t^{2}}{2}\|\tilde g_t\|^{2}.
\tag{4}
\]

\noindent\textbf{[Step 3] (Decompose the inner product).}
Add and subtract the unbiased stochastic gradient \(g_t\):
\[
\langle\nabla f(w_t),\tilde g_t\rangle
= \langle\nabla f(w_t),g_t\rangle
+ \langle\nabla f(w_t),\tilde g_t-g_t\rangle .
\tag{5}
\]

\noindent\textbf{[Step 4] (Take conditional expectation).}
Condition on \(\mathcal{F}_t:=\sigma(w_0,\dots,w_t)\).  Using
\(\E[g_t\mid\mathcal{F}_t]=\nabla f(w_t)\) (Assumption \ref{ass:unbiased}) we obtain
\[
\E\!\big[\langle\nabla f(w_t),g_t\rangle\mid\mathcal{F}_t\big]
= \|\nabla f(w_t)\|^{2}.
\tag{6}
\]

\noindent\textbf{[Step 5] (Bound the compression error term).}
Since \(\tilde g_t-g_t = \mathcal{C}_K(g_t)-g_t\) and
\(\E[\mathcal{C}_K(g_t)-g_t\mid\mathcal{F}_t]=0\) (the Top‑\(K\) operator is
unbiased in expectation over the random tie‑breaking), we have
\[
\E\!\big[\langle\nabla f(w_t),\tilde g_t-g_t\rangle\mid\mathcal{F}_t\big]=0.
\tag{7}
\]

\noindent\textbf{[Step 6] (Bound the squared norm of the compressed gradient).}
Write
\[
\|\tilde g_t\|^{2}
= \|g_t+(\tilde g_t-g_t)\|^{2}
= \|g_t\|^{2}+2\langle g_t,\tilde g_t-g_t\rangle+\|\tilde g_t-g_t\|^{2}.
\]
Taking conditional expectation and using that
\(\E[\tilde g_t-g_t\mid\mathcal{F}_t]=0\) yields
\[
\E\!\big[\|\tilde g_t\|^{2}\mid\mathcal{F}_t\big]
= \E\!\big[\|g_t\|^{2}\mid\mathcal{F}_t\big]
+ \E\!\big[\|\tilde g_t-g_t\|^{2}\mid\mathcal{F}_t\big].
\tag{8}
\]
Apply Lemma~\ref{lem:comp} to the second term:
\[
\E\!\big[\|\tilde g_t-g_t\|^{2}\mid\mathcal{F}_t\big]
\le (1-\rho)\,\E\!\big[\|g_t\|^{2}\mid\mathcal{F}_t\big].
\tag{9}
\]
Hence
\[
\E\!\big[\|\tilde g_t\|^{2}\mid\mathcal{F}_t\big]
\le \big(1+(1-\rho)\big)\,\E\!\big[\|g_t\|^{2}\mid\mathcal{F}_t\big]
\le \big(1+(1-\rho)\big)G^{2},
\tag{10}
\]
where the last inequality uses Assumption \ref{ass:unbiased}.

\noindent\textbf{[Step 7] (Take full expectation of \eqref{4}).}
Combine \eqref{4} with the results of Steps~4–6 and then take expectation:
\[
\begin{aligned}
\E\big[f(w_{t+1})\big]
&\le \E\big[f(w_t)\big]
-\eta_t \E\!\big[\|\nabla f(w_t)\|^{2}\big]
+\frac{L\eta_t^{2}}{2}\,
\E\!\big[\|\tilde g_t\|^{2}\big] .
\end{aligned}
\tag{11}
\]

\noindent\textbf{[Step 8] (Insert the bound (10) into (11)).}
\[
\E\big[f(w_{t+1})\big]
\le \E\big[f(w_t)\big]
-\eta_t \E\!\big[\|\nabla f(w_t)\|^{2}\big]
+\frac{L\eta_t^{2}}{2}\,
\big(1+(1-\rho)\big)G^{2}.
\tag{12}
\]

\noindent\textbf{[Step 9] (Rearrange).}
Move the gradient term to the left‑hand side:
\[
\eta_t \E\!\big[\|\nabla f(w_t)\|^{2}\big]
\le \E\big[f(w_t)-f(w_{t+1})\big]
+\frac{L\eta_t^{2}}{2}\,
\big(1+(1-\rho)\big)G^{2}.
\tag{13}
\]

\noindent\textbf{[Step 10] (Sum over \(t=0,\dots,T-1\)).}
Summing \eqref{13} yields a telescoping series:
\[
\sum_{t=0}^{T-1}\eta_t \E\!\big[\|\nabla f(w_t)\|^{2}\big]
\le f(w_0)-\E\big[f(w_T)\big]
+\frac{L G^{2}\big(1+(1-\rho)\big)}{2}
\sum_{t=0}^{T-1}\eta_t^{2}.
\tag{14}
\]

\noindent\textbf{[Step 11] (Lower bound the left‑hand side).}
Since \(\eta_t\) is non‑increasing, \(\eta_t\ge\eta_{T-1}\) for all
\(t\le T-1\).  Hence
\[
\eta_{T-1}\sum_{t=0}^{T-1}\E\!\big[\|\nabla f(w_t)\|^{2}\big]
\le \sum_{t=0}^{T-1}\eta_t \E\!\big[\|\nabla f(w_t)\|^{2}\big].
\tag{15}
\]

\noindent\textbf{[Step 12] (Combine (14) and (15)).}
\[
\eta_{T-1}\sum_{t=0}^{T-1}\E\!\big[\|\nabla f(w_t)\|^{2}\big]
\le f(w_0)-f^\star
+\frac{L G^{2}\big(1+(1-\rho)\big)}{2}
\sum_{t=0}^{T-1}\eta_t^{2},
\tag{16}
\]
where \(f^\star:=\inf_w f(w)\le \E[f(w_T)]\).

\noindent\textbf{[Step 13] (Insert the concrete step‑size
\(\eta_t=\frac{\eta}{\sqrt{t+1}}\)).}
Then
\[
\eta_{T-1}= \frac{\eta}{\sqrt{T}},\qquad
\sum_{t=0}^{T-1}\eta_t^{2}
= \eta^{2}\sum_{t=0}^{T-1}\frac{1}{t+1}
\le \eta^{2}\big(1+\log T\big)
\le 2\eta^{2}\log T,
\]
the last inequality holding for \(T\ge e\).  Plugging into \eqref{16} gives
\[
\frac{\eta}{\sqrt{T}}\sum_{t=0}^{T-1}\E\!\big[\|\nabla f(w_t)\|^{2}\big]
\le f(w_0)-f^\star
+\frac{L G^{2}\big(1+(1-\rho)\big)}{2}\,2\eta^{2}\log T.
\tag{17}
\]

\noindent\textbf{[Step 14] (Divide by \(T\) and simplify).}
Define the averaged gradient norm \(\Phi_T\) as in the theorem.
From \eqref{17}:
\[
\Phi_T
= \frac{1}{T}\sum_{t=0}^{T-1}\E\!\big[\|\nabla f(w_t)\|^{2}\big]
\le
\frac{\sqrt{T}}{\eta T}\big(f(w_0)-f^\star\big)
+\frac{L G^{2}\big(1+(1-\rho)\big)}{\rho}\,
\frac{\eta\log T}{\sqrt{T}}.
\]
Dropping the slowly growing \(\log T\) term (or absorbing it into the
constant by choosing \(\eta\) appropriately) yields the cleaner bound
\eqref{2}:
\[
\Phi_T \le
\frac{2\big(f(w_0)-f^\star\big)}{\eta\sqrt{T}}
\;+\;
\frac{L G^{2}\big(1+(1-\rho)\big)}{\rho}\,
\frac{\eta}{\sqrt{T}}.
\]
This completes the proof of Theorem \ref{thm:main}. \hfill\(\square\)

\subsection*{Rate Derivation (With Explicit Constants)}
Choosing
\[
\eta = \sqrt{\frac{2\big(f(w_0)-f^\star\big)\,\rho}
{L G^{2}\big(1+(1-\rho)\big)}}
\]
balances the two terms in \eqref{2} and gives
\[
\Phi_T \;\le\;
2\sqrt{\frac{L G^{2}\big(1+(1-\rho)\big)}{\rho}}
\;\frac{\sqrt{f(w_0)-f^\star}}{\sqrt{T}}.
\]
Thus the constant factor scales as \(\sqrt{(1-\rho)/\rho}\).

\subsection*{Special Cases}
\begin{itemize}
\item \textbf{Full communication (\(\rho=1\)).}
The bound reduces to
\(\Phi_T\le \frac{2(f(w_0)-f^\star)}{\eta\sqrt{T}}+
\frac{L G^{2}\eta}{\sqrt{T}}\), the classical SGD result.
\item \textbf{Extreme sparsity (\(\rho\ll1\)).}
The constant blows up as \(\mathcal{O}(\rho^{-1/2})\); nevertheless the
asymptotic rate \(\mathcal{O}(1/\sqrt{T})\) is preserved.
\end{itemize}

%%%%%%%%%%%%%%%%%%%%%%%%%%%%%%%%%%%%%%%%%%%%%%%%%%%%%%%%%%%%
\begin{thebibliography}{9}
\bibitem{Alistarh2017}
D.~Alistarh, J.~Li, R.~Tomioka, and M.~Vojnovic,
\newblock \emph{QSGD: Randomized Quantization for Communication-Efficient
Stochastic Gradient Descent},
\newblock in Proc. NIPS, 2017.
\end{thebibliography}

\end{document}